\documentclass[10pt,letterpaper]{article}
\usepackage[T1]{fontenc}
\usepackage[utf8]{inputenc}
\usepackage[spanish,es-tabla]{babel}
\usepackage{csquotes}
\usepackage{mathpazo}
\usepackage[scaled=0.9]{helvet}
\usepackage{courier}
\usepackage{calc}
\usepackage[top=2cm, bottom=2cm, left=2cm, right=2cm]{geometry}
\usepackage{xcolor}
\usepackage{booktabs}
\usepackage{longtable}
\usepackage{array}
\usepackage{ragged2e}
\usepackage{xurl}
\usepackage[strings]{underscore}
\usepackage{fancyhdr}
\usepackage{sectsty}
\usepackage{tocloft}
\usepackage[colorlinks=true,linkcolor=blue,urlcolor=blue,citecolor=blue,breaklinks=true]{hyperref}
\usepackage{enumitem}

\definecolor{uncpblue}{HTML}{1A237E}
\allsectionsfont{\color{uncpblue}}

\renewcommand{\cftsecleader}{\cftdotfill{\cftsecdotsep}}
\cftpagenumbersoff{section}
\cftpagenumbersoff{subsection}
\cftpagenumbersoff{subsubsection}
\renewcommand{\cftsecfont}{\normalfont}
\renewcommand{\cftsubsecfont}{\normalfont}
\renewcommand{\cftsubsubsecfont}{\normalfont}
\setlength{\parindent}{0pt}
\setlength{\parskip}{4pt}
\setlength{\tabcolsep}{3pt}
\renewcommand{\arraystretch}{0.85}

\pagestyle{fancy}
\fancyhf{}
\fancyhead[L]{\small\textcolor{uncpblue}{\textbf{SGSI-UNCP}}}
\fancyhead[R]{\small\textcolor{uncpblue}{F-SGSI-04}}
\fancyfoot[C]{\thepage}
\renewcommand{\headrulewidth}{0.4pt}
\renewcommand{\footrulewidth}{0.4pt}
\setlength{\headheight}{14pt}

\setlist{nosep,leftmargin=*,after=\vspace{2pt}}
\setlength{\emergencystretch}{2em}

\newcommand{\makecover}{
\begin{titlepage}
\centering
\vspace*{2cm}
{\Huge\bfseries\color{uncpblue} SGSI-UNCP\par}
\vspace{0.7cm}
{\LARGE Plantilla de Informe de No Conformidad y Acción Correctiva
(RAC)\par}
\vspace{0.5cm}
{\large Documento Base del Sistema de Gestión\\de Seguridad de la Información\par}
\vspace{1.5cm}
{\small Universidad Nacional del Centro del Perú\par}
\vfill
\end{titlepage}
}

\begin{document}

\makecover
\setcounter{tocdepth}{2}
\RaggedRight
\tableofcontents
\newpage

\RaggedRight
\section{Plantilla de Informe de No Conformidad y Acción Correctiva
(RAC)}\label{plantilla-de-informe-de-no-conformidad-y-acciuxf3n-correctiva-rac}

\textbf{Organización:} Universidad Nacional del Centro del Perú (UNCP)
\textbf{ID de RAC:} RAC-2026-001 \textbf{ID de Auditoría (si aplica):}
AUD-2026-001

\begin{center}\rule{0.5\linewidth}{0.5pt}\end{center}

\subsection{Descripción de la No
Conformidad}\label{descripciuxf3n-de-la-no-conformidad}

\textbf{Fecha de Detección:} 22/05/2026 \textbf{Detectado por:} Mg.
Carlos Ramos Quispe (Auditor Líder) \textbf{Fuente:} Auditoría Interna
AUD-2026-001

\textbf{Descripción Detallada:}

En la inspección a la sala de servidores de la Facultad de Ingeniería de
Sistemas se encontraron 3 equipos (2 switches Cisco Catalyst 2960, 1
servidor HP ProLiant DL380) no registrados en el inventario de activos
R-SGSI-01. No tienen código de activo ni custodio asignado. Esta
situación incumple el control A.7.8 de ISO 27001:2022 (``Los equipos
deben ser inventariados'') y la política POL-SGSI-07.

\begin{center}\rule{0.5\linewidth}{0.5pt}\end{center}

\subsection{Análisis de Causa Raíz}\label{anuxe1lisis-de-causa-rauxedz}

\textbf{Método Utilizado:} 5 Porqués \textbf{Resultado del Análisis:}

\begin{enumerate}
\def\labelenumi{\arabic{enumi}.}
\tightlist
\item
  ¿Por qué los equipos no están inventariados? Porque no se incluyeron
  en el levantamiento inicial de activos.
\item
  ¿Por qué no se incluyeron? Porque la sala de servidores de la FIIS no
  fue considerada en el alcance del inventario.
\item
  ¿Por qué no fue considerada? Porque el procedimiento R-SGSI-01 no
  especificaba claramente las ubicaciones a cubrir.
\item
  ¿Por qué no se especificaron? Porque el responsable del inventario
  asumió que solo el Data Center principal alberga servidores.
\item
  ¿Por qué asumió eso? Porque no existía una lista completa de
  ubicaciones con infraestructura TI ni un procedimiento de verificación
  cruzada.
\end{enumerate}

\textbf{Causa Raíz:} Ausencia de un censo completo de ubicaciones con
infraestructura TI al momento de elaborar el inventario, y falta de
verificación cruzada con registros de compras y redes.

\begin{center}\rule{0.5\linewidth}{0.5pt}\end{center}

\subsection{Plan de Acción (Corrección y Acción
Correctiva)}\label{plan-de-acciuxf3n-correcciuxf3n-y-acciuxf3n-correctiva}

\begin{itemize}
\tightlist
\item
  \textbf{Corrección (Acción inmediata):} Registrar los 3 equipos
  encontrados en el inventario R-SGSI-01 con código de activo, custodio
  asignado (Jefe de la FIIS) y clasificación correspondiente. Realizar
  etiquetado físico de los equipos.
\item
  \textbf{Acción Correctiva (Para evitar recurrencia):} Realizar un
  censo físico completo de todas las salas de servidores y gabinetes de
  telecomunicaciones en todas las facultades y dependencias de la UNCP.
  Actualizar el procedimiento R-SGSI-01 para incluir la verificación
  cruzada con registros de compras (DGA), redes (OTI) y planos de
  infraestructura.
\end{itemize}

\textbf{Responsable:} Ing. Marco Antonio Quispe Laura (Jefe de OTI)
\textbf{Fecha Límite:} 30/06/2026

\begin{center}\rule{0.5\linewidth}{0.5pt}\end{center}

\subsection{Verificación de la
Eficacia}\label{verificaciuxf3n-de-la-eficacia}

\textbf{Fecha de Verificación:} 15/07/2026 \textbf{Verificado por:} Ing.
Luis Castillo Gutierrez (Oficial de Seguridad)

\textbf{Resultado:}

\begin{itemize}
\tightlist
\item[$\boxtimes$]
  Eficaz (Se cierra el RAC)
\item[$\square$]
  No Eficaz (Se requiere nuevo análisis de causa)
\end{itemize}

\textbf{Evidencia de Verificación:}

Se realizó una inspección conjunta OTI-Oficial de Seguridad el
10/07/2026. Los 3 equipos están registrados en R-SGSI-01 con códigos
UNCP-SRV-025, UNCP-SW-047 y UNCP-SW-048. Se verificó el etiquetado
físico en la sala FIIS. Adicionalmente, se completó el censo de 12
ubicaciones adicionales (facultades), identificando 5 equipos más que
fueron incorporados al inventario. El procedimiento R-SGSI-01 fue
actualizado (versión 2.1) incluyendo el paso de verificación cruzada con
DGA y OTI Redes.

\begin{center}\rule{0.5\linewidth}{0.5pt}\end{center}

\textbf{Firmas:}

\begin{center}\rule{0.5\linewidth}{0.5pt}\end{center}

\textbf{Ing. Marco Antonio Quispe Laura} Responsable del Área (OTI)

\begin{center}\rule{0.5\linewidth}{0.5pt}\end{center}

\textbf{Ing. Luis Castillo Gutierrez} Oficial de Seguridad y Confianza
Digital

\end{document}
